\paperchapter{December 2024}{1 December 2024}{A}{50}{18th GSET}
\keytable{\kcell{1}{A} & \kcell{2}{C} & \kcell{3}{A} & \kcell{4}{D} & \kcell{5}{A} & \kcell{6}{B} & \kcell{7}{A} & \kcell{8}{B} & \kcell{9}{A} & \kcell{10}{A} \\
\kcell{11}{A} & \kcell{12}{D} & \kcell{13}{C} & \kcell{14}{D} & \kcell{15}{D} & \kcell{16}{B} & \kcell{17}{B} & \kcell{18}{B} & \kcell{19}{D} & \kcell{20}{B} \\
\kcell{21}{A} & \kcell{22}{B} & \kcell{23}{B} & \kcell{24}{B} & \kcell{25}{A} & \kcell{26}{D} & \kcell{27}{A} & \kcell{28}{A} & \kcell{29}{D} & \kcell{30}{A} \\
\kcell{31}{C} & \kcell{32}{A} & \kcell{33}{B} & \kcell{34}{A} & \kcell{35}{B} & \kcell{36}{A} & \kcell{37}{A} & \kcell{38}{C} & \kcell{39}{B} & \kcell{40}{D} \\
\kcell{41}{B} & \kcell{42}{C} & \kcell{43}{C} & \kcell{44}{B} & \kcell{45}{B} & \kcell{46}{C} & \kcell{47}{B} & \kcell{48}{D} & \kcell{49}{B} & \kcell{50}{C}}
\begin{question}{1}
{Semantic barrier in communication is\par}
\optstwo{Use of technical/difficult language}{Use of technology}{Lack of attention}{Lack of co-ordination}
\answer{A}{Use of technical/difficult language}
\begin{explain}
A \textbf{semantic barrier} arises from problems of \emph{meaning}: words, symbols, jargon or technical terms that the receiver does not understand, or understands differently. Using technical or difficult language is the classic example.
\why{Use of technology relates to channel or technological barriers; lack of attention is a psychological barrier; lack of co-ordination is an organisational barrier.}
\end{explain}
\end{question}

\begin{question}{2}
{The key purpose of effective communication is\par}
\optstwo{Create barriers}{Promote misunderstanding}{Effect change}{Confuse audience}
\answer{C}{Effect change}
\begin{explain}
Communication is purposeful: it aims to inform, persuade, motivate and so \textbf{bring about a change} in the knowledge, attitude or behaviour of the receiver. Effective communication achieves the intended effect.
\why{Barriers, misunderstanding and confusion are the \emph{opposite} of effective communication.}
\end{explain}
\end{question}

\begin{question}{3}
{Identify the correct sequence of the elements in the communication system.\newline a. Speech\newline b. Speaker\newline c. Receiver\newline d. Effect\newline The correct choice is :\par}
\optsfour{b, a, c, d}{a, b, d, c}{b, a, d, c}{c, d, b, a}
\answer{A}{b, a, c, d}
\begin{explain}
In \textbf{Aristotle's model} of communication (the earliest linear model) the elements occur in this order:
\begin{center}
\begin{tikzpicture}[node distance=6mm,every node/.style={draw=brand,rounded corners,fill=soft,font=\small,minimum height=6mm}]
\node (a) {Speaker (b)}; \node[right=of a] (b) {Speech (a)}; \node[right=of b] (c) {Receiver (c)}; \node[right=of c] (d) {Effect (d)};
\draw[-Stealth,thick,brand] (a)--(b); \draw[-Stealth,thick,brand] (b)--(c); \draw[-Stealth,thick,brand] (c)--(d);
\end{tikzpicture}
\end{center}
Speaker $\rightarrow$ Speech $\rightarrow$ Audience (receiver) $\rightarrow$ Effect, i.e.\ \textbf{b, a, c, d}. (Lasswell's later model follows the same logic: \emph{Who says What, in Which channel, to Whom, with What effect}.)
\end{explain}
\end{question}

\begin{question}{4}
{Which of the following is an example of non-verbal communication ?\par}
\optstwo{Writing on the blackboard}{Speaking on microphone}{Sending a voice message}{Nodding head to indicate disagreement}
\answer{D}{Nodding head to indicate disagreement}
\begin{explain}
\textbf{Non-verbal} communication uses no words: gestures, facial expressions, posture, eye contact, head movements, touch, space and time. \textbf{Nodding the head} (to agree or disagree) is a gesture, which falls under kinesics.
\why{Writing on the blackboard is written verbal communication; speaking on a microphone and sending a voice message are oral verbal communication.}
\end{explain}
\end{question}

\begin{question}{5}
{'Encoding' in the communication system means\par}
\optsone{Converting thoughts or ideas into a message}{Translating messages into ideas or thoughts}{Providing opinions on the sender}{Avoiding difficult words in the message}
\answer{A}{Converting thoughts or ideas into a message}
\begin{explain}
\textbf{Encoding} is the sender's act of converting thoughts, ideas or feelings into a message using symbols (words, pictures, gestures).
\why{Option (B) describes \emph{decoding}, which is done by the receiver.}
\end{explain}
\end{question}

\begin{question}{6}
{Which one among the following is the common language for AI ?\par}
\optsfour{Lisp}{Python}{PHP}{JAVA}
\answer{B}{Python}
\begin{explain}
\textbf{Python} is the most widely used language in present-day AI and machine learning, thanks to libraries such as TensorFlow, PyTorch, scikit-learn, NumPy and Keras. The key treats it as the common language for AI.
\why{LISP (1958) was the historical language of AI research, but today it is niche. PHP is a web scripting language, and Java, though usable, is not the common AI language.}
\end{explain}
\end{question}

\begin{question}{7}
{Regarding ICT based governance, which of the following statements is/are correct?\newline i. Increasing transparency is an advantage of ICT based governance.\newline ii. e-governance is related to automation of administrative processes.\newline iii. Due to ICT based governance, workforce gets reduced.\newline Choose the correct option:\par}
\optstwo{i, ii and iii}{i and ii only}{i and iii only}{ii and iii only}
\answer{A}{i, ii and iii}
\begin{explain}
All three statements are correct:
\begin{itemize}
\item ICT increases \textbf{transparency} (online records, RTI portals, e-tendering).
\item e-Governance involves \textbf{automation of administrative processes} (e-office, online services).
\item Automation reduces paperwork and manual jobs, so the \textbf{workforce needed is reduced}.
\end{itemize}
Hence i, ii and iii.
\end{explain}
\end{question}

\begin{question}{8}
{Which among the following are the National Co-ordinators for SWAYAM?\newline 1. Indian Institute of Technology Kanpur (IITK)\newline 2. University Grants Commission (UGC)\newline 3. Consortium for Educational Communication (CEC)\newline 4. Indian Institute of Management Bangalore (IIMB)\newline 5. National Institute of Open Schooling (NIOS)\newline Choose the correct answer from the options given below :\par}
\optsfour{1, 2, 3 and 4}{2, 3, 4 and 5}{1, 2, 3 and 5}{2, 3 and 4}
\answer{B}{2, 3, 4 and 5}
\begin{explain}
SWAYAM has \textbf{nine National Coordinators}, each responsible for one type of content: AICTE (self-paced and international courses), \textbf{NPTEL/IIT Madras} (engineering), \textbf{UGC} (non-technical postgraduate education), \textbf{CEC} (undergraduate education), \textbf{NCERT} and \textbf{NIOS} (school education), IGNOU (out-of-school students), \textbf{IIM Bangalore} (management studies) and NITTTR (teacher training).
Among the given names, UGC, CEC, IIMB and NIOS are national coordinators. IIT Kanpur is not (NPTEL is coordinated through IIT Madras), so the answer is \textbf{2, 3, 4 and 5}.
\end{explain}
\end{question}

\begin{question}{9}
{Given below are two statements:\newline Statement I: SWAYAM is India's own MOOC platform, which offers free online courses in almost all the discipline to students/learners.\newline Statement - II: Coursera, Udemy, Udacity, Iversity are other MOOC platforms.\newline In light of the above statements, choose the correct answer from the options given below\par}
\optsone{Both statement I and statement II are true}{Both statement I and statement II are false}{Statement I is true but statement II is false}{Statement I is false but statement II is true}
\answer{A}{Both statement I and statement II are true}
\begin{explain}
\textbf{Statement I} is true: SWAYAM (Study Webs of Active-Learning for Young Aspiring Minds), launched in 2017 by MHRD, is India's own MOOC platform offering free courses from Class~9 to postgraduate level. \textbf{Statement II} is true: Coursera, Udemy, Udacity and iversity are well-known international MOOC providers.
\end{explain}
\end{question}

\begin{question}{10}
{Given below are two statements, one is labeled as Assertion (A) and the other is labeled as Reason (R).\newline Assertion (A): 'Touch screen' is considered as both input and output device.\newline Reason (R): 'Touch screen' can send data and receive data at the same time.\newline In light of the above statements, choose the correct answer from the options given below :\par}
\optsone{Both (A) and (R) are true and (R) is the correct explanation of (A)}{Both (A) and (R) are true but (R) is not the correct explanation of (A)}{(A) is true but (R) is false}{(A) is false but (R) is true}
\answer{A}{Both (A) and (R) are true and (R) is the correct explanation of (A)}
\begin{explain}
A touch screen both \textbf{displays} output and \textbf{accepts} touch input, so it is an input as well as an output device. The reason, that it can send and receive data at the same time, correctly explains why it plays both roles. Both are true and (R) explains (A).
\end{explain}
\end{question}

\begin{question}{11}
{Which of the following components fall under presage stage of presage processor-product model of teaching-learning process ?\newline 1. Content analysis\newline 2. Learning testing\newline 3. Classroom interaction\newline 4. Learning outcome plan\newline 5. Learning supervision\newline 6. Learner's motivation\newline The correct choice is :\par}
\optsfour{1, 4, 6}{2, 3, 5}{1, 5, 6}{1, 4, 5}
\answer{A}{1, 4, 6}
\begin{explain}
In Dunkin and Biddle's \textbf{Presage--Process--Product} model:
\begin{itemize}
\item \textbf{Presage} variables come \emph{before} teaching: teacher characteristics, planning, \emph{content analysis}, the \emph{learning outcome plan}, and learner characteristics such as \emph{motivation}.
\item \textbf{Process} variables operate \emph{during} teaching: classroom interaction and learning supervision.
\item \textbf{Product} variables are the \emph{results}: learning and its testing.
\end{itemize}
Presage components here are 1, 4 and 6.
\end{explain}
\end{question}

\begin{question}{12}
{Discussion method of teaching does NOT have the advantage of developing\par}
\optsone{Understanding and experiencing}{Reflective level of experiencing}{Active learning}{Organisation of content flow in a structured sequence}
\answer{D}{Organisation of content flow in a structured sequence}
\begin{explain}
The discussion method promotes understanding, reflective thinking and active participation, but because it follows the flow of the group's ideas it does \textbf{not} guarantee that content is organised in a structured, logical sequence. That is the strength of the lecture or programmed method.
\end{explain}
\end{question}

\begin{question}{13}
{A teacher through experience found that learners were able to memorise and recall content, if given in brief summarised form. Of the given below which one is NOT the implication ?\par}
\optsone{Brief summarisation will motivate learner for active participation}{Brief summarisation will improve understanding and conceptualisation}{Brief summarisation will focus on memorisation irrespective of understanding}{Brief summarisation is key for cognitive learning}
\answer{C}{Brief summarisation will focus on memorisation irrespective of understanding}
\begin{explain}
The teacher observed that brief summaries help learners \emph{memorise and recall}. Sensible implications are that summaries motivate participation, aid understanding and conceptualisation, and support cognitive learning. Saying that summarisation focuses on memorisation \textbf{irrespective of understanding} contradicts good pedagogy and is \emph{not} a valid implication.
\end{explain}
\end{question}

\begin{question}{14}
{The online teaching has advantage with reference to\par}
\optstwo{use of multimedia system}{well structuring of teaching}{objective based teaching}{all of the above}
\answer{D}{all of the above}
\begin{explain}
Online teaching allows rich \textbf{multimedia} (video, animation, simulation), requires content to be \textbf{well structured} into modules, and is built around clear \textbf{learning objectives}. All three are advantages, so the answer is \emph{all of the above}.
\end{explain}
\end{question}

\begin{question}{15}
{Which of the following learner characteristics disturb classroom teaching?\newline 1. Self confidence\newline 2. Absence of teacher clarification\newline 3. Logical arguments for clarification\newline 4. Classroom noise by learners\newline 5. Peer attention seeking argument\newline 6. Access to large amount of digital material on net\newline Choose the correct option:\par}
\optsfour{1, 2, 5}{2, 4, 6}{3, 4, 6}{2, 4, 5}
\answer{D}{2, 4, 5}
\begin{explain}
Factors that \emph{disturb} classroom teaching:
\begin{itemize}
\item 2. absence of teacher clarification, which leaves learners confused;
\item 4. classroom noise by learners;
\item 5. peer attention-seeking arguments.
\end{itemize}
Self-confidence (1) and logical arguments for clarification (3) help learning, and access to digital material (6) is a resource rather than a disturbing \emph{characteristic}. The answer is 2, 4, 5.
\end{explain}
\end{question}

\begin{question}{16}
{In the following series of numbers, one number is wrong in the series :\newline 864, 420, 200, 96, 40, 16, 6\newline Which is the wrong number?\par}
\optsfour{40}{96}{200}{420}
\answer{B}{96}
\begin{explain}
Rule: each term $=\dfrac{\text{previous}}{2}-k$, with $k$ decreasing by 2:
\[864\xrightarrow{\div2-12}420\xrightarrow{\div2-10}200\xrightarrow{\div2-8}\mathbf{92}\xrightarrow{\div2-6}40\xrightarrow{\div2-4}16\xrightarrow{\div2-2}6\]
The series should have \textbf{92}, so \textbf{96} is the wrong number (and $92/2-6=40$ confirms it).
\end{explain}
\end{question}

\begin{question}{17}
{The difference between compound and simple interest on a certain sum of money at the interest rate of 10\% per annum for 1.5 years is 183, when the interest is compounded semi annually. What is the sum of money?\par}
\optsfour{19,200}{24,000}{35,000}{38,060}
\answer{B}{24,000}
\begin{explain}
Semi-annual: rate $=5\%$ per half-year, $n=3$ half-years.
\[CI=P[(1.05)^3-1]=0.157625P,\qquad SI=P\times10\%\times1.5=0.15P\]
\[CI-SI=0.007625P=183\ \Rightarrow\ P=\frac{183}{0.007625}=\textbf{Rs.\ 24{,}000}\]
\end{explain}
\end{question}

\begin{question}{18}
{The following table gives the percentage distribution of population of five States A, B, C, D and E on the basis of poverty line and also on the basis of gender:\par}
\qfigure{\small
\begin{tabular}{c c cc}\toprule
& & \multicolumn{2}{c}{\textbf{Males : Females}}\\\cmidrule(l){3-4}
\textbf{State} & \textbf{\% below poverty line} & \textbf{Below PL} & \textbf{Above PL}\\\midrule
A & 35 & 5 : 6 & 6 : 7\\ B & 25 & 3 : 5 & 4 : 5\\ C & 24 & 1 : 2 & 2 : 3\\ D & 19 & 3 : 2 & 4 : 3\\ E & 15 & 5 : 3 & 3 : 2\\\bottomrule
\end{tabular}}
What will be the number of females above poverty line in the State D if it is known that the population of State D is 7 million?
\optsfour{2.87 million}{2.43 million}{2.17 million}{1.98 million}
\answer{B}{2.43 million}
\begin{explain}
State D: 19\% below poverty line, so $100-19=81\%$ are above.\\
Population above poverty line $=81\%\times7=5.67$ million.\\
Males : Females above PL $=4:3$, so females $=\tfrac37\times5.67=\textbf{2.43 million}$.
\end{explain}
\end{question}

\begin{question}{19}
{The information about the students of a college is given below :\newline Members of Cricket Club only: 5\%\newline Members of Football Club only: 15\%\newline Members of both the Club: 75\%\newline Non members: 40\newline What is the number of students who are members of both the Clubs ?\par}
\optsfour{450}{500}{550}{600}
\answer{D}{600}
\begin{explain}
Cricket only $+$ Football only $+$ Both $=5+15+75=95\%$, so non-members $=5\%$.
\[5\%\text{ of total}=40\ \Rightarrow\ \text{total}=800\qquad \text{Both clubs}=75\%\times800=\mathbf{600}\]
\end{explain}
\end{question}

\begin{question}{20}
{The number of auto-parts manufactured by 5 identical machines during 5 days is given below:\par}
\qfigure{\small
\begin{tabular}{lccccc}\toprule
\textbf{Day} $\backslash$ \textbf{Machine} & \textbf{A} & \textbf{B} & \textbf{C} & \textbf{D} & \textbf{E}\\\midrule
Monday & 139 & 147 & 211 & 141 & 184\\ Tuesday & 141 & 189 & 164 & 189 & 151\\ Wednesday & 115 & 141 & 159 & 156 & 136\\ Thursday & 89 & 223 & 120 & 147 & 113\\ Friday & 187 & 93 & 257 & 160 & 124\\\bottomrule
\end{tabular}}
The number of auto-parts manufactured by Machine E on Monday is what percent more than the number of parts manufactured by Machine C on Thursday?
\optsfour{47.66\%}{53.33\%}{59.66\%}{62.33\%}
\answer{B}{53.33\%}
\begin{explain}
Machine E on Monday $=184$; Machine C on Thursday $=120$.
\[\frac{184-120}{120}\times100=\frac{64}{120}\times100=\mathbf{53.33\%}\]
\end{explain}
\end{question}

\begin{question}{21}
{Which agency in India conducts NEET Examination for admission into medical colleges ?\par}
\optstwo{National Testing Agency}{National Board of Assessment}{Medical Council of India}{All India Institute of Medical Science}
\answer{A}{National Testing Agency}
\begin{explain}
Since 2019 \textbf{NEET-UG} has been conducted by the \textbf{National Testing Agency (NTA)}, an autonomous body set up in 2017 under the Ministry of Education. It also conducts JEE (Main), CUET and UGC-NET.
\why{CBSE conducted NEET before 2019. The Medical Council of India was replaced by the National Medical Commission in 2020.}
\end{explain}
\end{question}

\begin{question}{22}
{'ICSSR' stands for\par}
\optsone{International Council for Social Science Research}{Indian Council of Social Science Research}{Indian Council of Scientific Sustainable Research}{International Conference on Solid Science Research}
\answer{B}{Indian Council of Social Science Research}
\begin{explain}
\textbf{ICSSR} is the \textbf{Indian Council of Social Science Research}, established in 1969 under the Ministry of Education to promote research in the social sciences.
\tip{Other councils: ICHR (Historical Research, 1972), ICPR (Philosophical Research, 1977), ICAR (Agricultural Research), ICMR (Medical Research), CSIR (Scientific \& Industrial Research).}
\end{explain}
\end{question}

\begin{question}{23}
{Which Institution is scored rank one (1), according to India Overall Rankings-2024 of NIRF, Govt. of India ?\par}
\optstwo{IISC, Bangalore}{IIT, Madras}{IIT, Delhi}{JNU, Delhi}
\answer{B}{IIT, Madras}
\begin{explain}
In \textbf{NIRF 2024} (India Rankings, Ministry of Education), \textbf{IIT Madras} was ranked No.~1 in the \emph{Overall} category for the sixth consecutive year (2019--2024). IISc Bangalore was 2nd and IIT Bombay 3rd.
\tip{In the \emph{University} category, IISc Bangalore is ranked 1st.}
\end{explain}
\end{question}

\begin{question}{24}
{'NIRF' stands for\par}
\optsone{National Industrial Ranking Framework}{National Institutional Ranking Framework}{National Individual Ranking Framework}{National Innovative Research Framework}
\answer{B}{National Institutional Ranking Framework}
\begin{explain}
\textbf{NIRF} is the \textbf{National Institutional Ranking Framework}, launched by MHRD in 2015 (first rankings 2016). Its five parameters are Teaching, Learning \& Resources; Research \& Professional Practice; Graduation Outcomes; Outreach \& Inclusivity; and Perception.
\end{explain}
\end{question}

\begin{question}{25}
{Name of the centres established to develop competencies among faculty members in Higher Education for better teaching, learning, research and academic leadership for effective implementation of NEP-2020\par}
\optsone{Malaviya Mission Teacher Training Centres}{Pandit Madan Mohan Malaviya Training Centres}{National Centre for Professional Development of Teacher}{Human Development and Research Centre}
\answer{A}{Malaviya Mission Teacher Training Centres}
\begin{explain}
The \textbf{Malaviya Mission Teacher Training Programme (MMTTP)} was launched by the Ministry of Education in 2023, renaming and revamping the earlier Human Resource Development Centres (formerly Academic Staff Colleges). Its \textbf{Malaviya Mission Teacher Training Centres} build the competencies of higher-education faculty in line with NEP~2020.
\end{explain}
\end{question}

\begin{question}{26}
{Which of the following is NOT a characteristic of basic research ?\par}
\optsone{Knowledge findings are replicable}{Generalisation of knowledge claims}{Concept and procedure of knowledge creation are publicly accessible}{Research knowledge is human centric}
\answer{D}{Research knowledge is human centric}
\begin{explain}
Basic (fundamental) research aims at \emph{general, objective} knowledge. Its findings are replicable, its claims generalisable, and its concepts and procedures publicly accessible for verification. Being \textbf{human-centric} (value-laden, focused on human problems) is not a defining characteristic of basic research; it describes applied or action-oriented research better.
\end{explain}
\end{question}

\begin{question}{27}
{Which one of the following is NOT a tool for quantitative data analysis in educational research?\par}
\optsfour{MS Word}{MS Excel}{SPSS}{Math Lab}
\answer{A}{MS Word}
\begin{explain}
\textbf{MS Word} is a word processor for creating documents; it is not a data-analysis tool.
\why{MS Excel (functions, pivot tables, Data Analysis ToolPak), SPSS (Statistical Package for the Social Sciences) and MATLAB are all used for quantitative analysis.}
\end{explain}
\end{question}

\begin{question}{28}
{A researcher intends to explore the results of possible factors for the organisation of 'Girl student support system development in the college campus'. Which research method will be most appropriate for the study ?\par}
\optstwo{Descriptive survey method}{Historical method}{Philosophical method}{Experimental method}
\answer{A}{Descriptive survey method}
\begin{explain}
The researcher wants to \emph{explore} the factors behind organising a support system for girl students as it exists or could exist on campus, typically by gathering views from students, teachers and administrators. Collecting such present-status information from a group is the \textbf{descriptive survey method}.
\why{Historical research studies past events; philosophical research analyses concepts and values; experimental research manipulates variables under control, which is not the aim here.}
\end{explain}
\end{question}

\begin{question}{29}
{What percentage of observation are broadly covered in $\pm2$ standard deviation (2$\sigma$) of a normal probability distribution ?\par}
\optsfour{60\%}{68\%}{90\%}{95\%}
\answer{D}{95\%}
\begin{explain}
In a normal distribution: $\pm1\sigma$ covers 68.26\%, $\pm2\sigma$ covers \textbf{95.44\%} ($\approx$ 95\%), and $\pm3\sigma$ covers 99.73\%.
\end{explain}
\end{question}

\begin{question}{30}
{Which one of the following is called non-probabilistic sampling?\par}
\optstwo{Quota sampling}{Cluster sampling}{Systematic sampling}{Random sampling}
\answer{A}{Quota sampling}
\begin{explain}
In \textbf{non-probability sampling} not every unit has a known chance of selection. Examples are quota, purposive, convenience, snowball and judgement sampling. \textbf{Quota sampling} is non-probabilistic.
\why{Cluster, systematic and simple random sampling are all \emph{probability} sampling techniques.}
\end{explain}
\end{question}

\begin{question}{31}
{Given : $U=-2$, $V=\frac{1}{2}$, $W=3$, $X=-4$, $Y=9$ and $Z=\frac{1}{6}$ . What is the value of $\sqrt{\frac{(W-2)^{2}}{V}+\frac{(Y-5)^{2}}{Z}+U-X}$ ?\par}
\optsfour{8}{9}{10}{11}
\answer{C}{10}
\begin{explain}
\[\frac{(W-2)^2}{V}=\frac{(3-2)^2}{1/2}=2,\qquad \frac{(Y-5)^2}{Z}=\frac{(9-5)^2}{1/6}=96,\qquad U-X=-2-(-4)=2\]
\[\sqrt{2+96+2}=\sqrt{100}=\mathbf{10}\]
\end{explain}
\end{question}

\begin{question}{32}
{If the compound interest on a certain sum at $16\frac{2}{3}\%$ for 3 years is 1,270, what is the simple interest on the same sum at the same rate and for the same period ?\par}
\optsfour{1,080}{2,010}{2,410}{2,560}
\answer{A}{1,080}
\begin{explain}
$16\tfrac23\%=\tfrac16$. The CI factor for 3 years is $\left(\tfrac76\right)^3-1=\tfrac{343}{216}-1=\tfrac{127}{216}$.
\[P=1270\times\frac{216}{127}=2160\qquad SI=2160\times\frac16\times3=\textbf{Rs.\ 1{,}080}\]
\end{explain}
\end{question}

\begin{question}{33}
{Complete the letter series:\newline ab\_a\_b\_aab\_aabb\_\par}
\optsfour{aabba}{babba}{babab}{ababa}
\answer{B}{babba}
\begin{explain}
Filling in \textbf{b, a, b, b, a} gives
\[\text{ab\underline{b}a}\ \ \text{\underline{a}b\underline{b}a}\ \ \text{ab\underline{b}a}\ \ \text{abb\underline{a}}\]
The block ``abba'' repeats four times, so the answer is ``babba'', option (B).
\end{explain}
\end{question}

\begin{question}{34}
{The average of three numbers is given as 44. If the second number is twice the first and is also thrice the third, what is the largest number ?\par}
\optsfour{72}{68}{64}{58}
\answer{A}{72}
\begin{explain}
Let the second number be $6k$. Then the first is $3k$ and the third is $2k$.
\[3k+6k+2k=3\times44=132\Rightarrow 11k=132\Rightarrow k=12\]
The numbers are 36, 72 and 24, so the largest is \textbf{72}.
\end{explain}
\end{question}

\begin{question}{35}
{10\% of the inhabitants of a village having died of cholera, during which 20\% remaining inhabitants left the village. The population is then reduced to 5040. What is the original inhabitants ?\par}
\optsfour{6000}{7000}{8000}{9000}
\answer{B}{7000}
\begin{explain}
After 10\% die, 90\% remain; then 20\% of these leave, so 80\% of 90\% remain.
\[P\times0.9\times0.8=5040\Rightarrow 0.72P=5040\Rightarrow P=\mathbf{7000}\]
\end{explain}
\end{question}

\begin{question}{36}
{Match the following and choose the correct option from given below:\par\smallskip\begin{tabularx}{\linewidth}{@{}>{\raggedright\arraybackslash}X>{\raggedright\arraybackslash}X@{}}\toprule
\textbf{SDG Number} & \textbf{Goal}\\\midrule
I. 4 & a. Clean water and sanitation\\
II. 6 & b. Climate action\\
III. 7 & c. Quality education\\
IV. 13 & d. Affordable and clean energy\\
\bottomrule\end{tabularx}\par\smallskip
\par}
\optstwo{I-c, II-a, III-d, IV-b}{I-c, II-d, III-a, IV-b}{I-a, II-b, III-c, IV-d}{I-b, II-c, III-a, IV-d}
\answer{A}{I-c, II-a, III-d, IV-b}
\begin{explain}
SDG~4 Quality education (c); SDG~6 Clean water \& sanitation (a); SDG~7 Affordable \& clean energy (d); SDG~13 Climate action (b). The match is \textbf{I-c, II-a, III-d, IV-b}.
\end{explain}
\end{question}

\begin{question}{37}
{Given below are two statements:\newline Statement I: The environmental disasters are natural and hence these are also termed as natural hazardous processes.\newline Statement -II: Earthquake, Tsunamis, Flood and Avalanche are examples of natural hazards.\newline In light of the above statements, choose the correct answer from the options given below :\par}
\optsone{Both statement I and statement II are true}{Both statement I and statement II are false}{Statement I is true but statement II is false}{Statement I is false but statement II is true}
\answer{A}{Both statement I and statement II are true}
\begin{explain}
The key treats both statements as true. Earthquakes, tsunamis, floods and avalanches are standard examples of \textbf{natural hazards} (Statement II). Statement I describes the disasters referred to as arising from natural hazardous processes, and the official key accepts it as true.
\tip{Hazard = a potentially damaging natural process. Disaster = a hazard that actually causes serious loss because the community is vulnerable.}
\end{explain}
\end{question}

\begin{question}{38}
{Which of the following are effect of global warming?\newline 1. Change in the rainfall pattern.\newline 2. Rise in the sea level.\newline 3. Loss of biodiversity.\newline 4. Shrinking of forest cover.\newline Choose the correct answer from the options given below :\par}
\optsfour{1, 2 and 3}{1, 3 and 4}{1, 2, 3 and 4}{1, 2 and 4}
\answer{C}{1, 2, 3 and 4}
\begin{explain}
Global warming leads to changed rainfall patterns, sea-level rise (thermal expansion and melting ice), loss of biodiversity, and shrinking forests (droughts, forest fires, species shifts). All four are effects, so the answer is (C).
\end{explain}
\end{question}

\begin{question}{39}
{Which of the following cyclone has affected the West Coast of India in last few years ?\newline 1. Cyclone Biparjoy\newline 2. Cyclone Fani\newline 3. Cyclone Tauktae\newline 4. Cyclone Amphan\newline Choose the correct answer from the options given below :\par}
\optsfour{1, 2 and 3}{1 and 3}{2 and 4}{1, 3 and 4}
\answer{B}{1 and 3}
\begin{explain}
West coast (Arabian Sea) cyclones: \textbf{Biparjoy} (June 2023, landfall near Jakhau, Kutch, Gujarat) and \textbf{Tauktae} (May 2021, landfall in Saurashtra, Gujarat).
East coast (Bay of Bengal) cyclones: \textbf{Fani} (2019, Odisha) and \textbf{Amphan} (2020, West Bengal).
The answer is 1 and 3.
\end{explain}
\end{question}

\begin{question}{40}
{The soil pollutant among the following are:\newline 1. Pesticides\newline 2. Radioactive wastes\newline 3. Municipal wastes\newline 4. Mining wastes\newline 5. Electronic wastes\newline Choose the correct answer from the options given below :\par}
\optstwo{1, 2, 3 and 5}{1, 2, 4 and 5}{1, 3, 4 and 5}{1, 2, 3, 4 and 5}
\answer{D}{1, 2, 3, 4 and 5}
\begin{explain}
Soil is polluted by pesticides and fertilisers, radioactive waste, municipal solid waste, mining waste (heavy metals, acid mine drainage) and e-waste (lead, mercury, cadmium). All five are soil pollutants.
\end{explain}
\end{question}

\begin{context}
\textbf{Directions for Q.~41 to 45:} Read the following paragraph and answer the questions.\par\smallskip
There are several other fields in which it is desirable that a clear policy should be laid down on an all-India basis. One of the fundamental rights guaranteed by the Constitution is the right of every citizen of the Union to free and compulsory education up to the age of 14. For the proper functioning of democracy, the Centre must see that every individual is equipped with the necessary knowledge, skill and aptitudes to discharge his duties as a responsible and co-operative citizen. What James Madison said about his country, the United States, many years ago, holds good today in our country also. ``A popular Government without popular information or the means of acquiring it is but a prologue to a farce of tragedy or perhaps both. Knowledge will forever govern ignorance, and the people who mean to be their own governors must arm themselves with the power which knowledge gives.'' We should, however like to add that it is not only knowledge that is required, but also the right kind of social training and the inculcation of right ideals without which knowledge by itself may be sterile or worse. Training for democracy postulates a balanced education in which social virtues, intellectual development and practical skill all receive due consideration and the pattern of such an education must be envisaged on an all-India basis.
\end{context}
\begin{question}{41}
{Which fundamental right is guaranteed by the Constitution of India as mentioned in the paragraph ?\par}
\optsone{The right to vote}{The right to free and compulsory education}{The right to freedom of speech}{The right to work and employment}
\answer{B}{The right to free and compulsory education}
\begin{explain}
The passage states: ``One of the fundamental rights guaranteed by the Constitution is the right of every citizen \ldots\ to \textbf{free and compulsory education} up to the age of 14.'' (Article 21A was added by the 86th Amendment, 2002, and implemented through the RTE Act, 2009.)
\end{explain}
\end{question}

\begin{question}{42}
{According to the paragraph, why is it important for every individual to acquire knowledge, skills and aptitudes ?\par}
\optsone{To become successful in life}{To participate in the global economy}{To discharge duties as a responsible and co-operative citizen}{To obtain government jobs}
\answer{C}{To discharge duties as a responsible and co-operative citizen}
\begin{explain}
``\ldots every individual is equipped with the necessary knowledge, skill and aptitudes \textbf{to discharge his duties as a responsible and co-operative citizen}.''
\end{explain}
\end{question}

\begin{question}{43}
{What did James Madison say about popular government ?\par}
\optsone{It cannot function without wealth.}{It must be led by knowledgeable people.}{Popular government without popular information is like a farce of tragedy.}{Democracy is the best form of government.}
\answer{C}{Popular government without popular information is like a farce of tragedy.}
\begin{explain}
Madison said that a popular government without popular information ``is but a \textbf{prologue to a farce of tragedy}.'' Option (C) paraphrases this.
\end{explain}
\end{question}

\begin{question}{44}
{What does 'training for democracy' postulate according to the paragraph ?\par}
\optsone{Training focused purely on intellectual development}{A balanced education that includes social virtues, intellectual development and practical skills}{A technical education system}{Education focused on leadership and politics}
\answer{B}{A balanced education that includes social virtues, intellectual development and practical skills}
\begin{explain}
``Training for democracy postulates a \textbf{balanced education} in which \textbf{social virtues, intellectual development and practical skill} all receive due consideration.''
\end{explain}
\end{question}

\begin{question}{45}
{Why does the author believe that educational policies should be established as an all-India basis ?\par}
\optsone{To promote political unity in India}{To ensure uniformity and balanced education across the country}{To provide higher education over primary education}{To ensure more students join government services}
\answer{B}{To ensure uniformity and balanced education across the country}
\begin{explain}
The author wants a policy on an all-India basis so that every citizen receives the same kind of \textbf{balanced education}, with uniform standards across the country. The pattern ``must be envisaged on an all-India basis.''
\end{explain}
\end{question}

\begin{question}{46}
{'X' is East of 'Y' which is in the North of 'Z'. If 'P' is to the East of 'Z', then in which direction of 'X' is 'P'?\par}
\optsfour{South-West}{South-East}{South}{West}
\answer{C}{South}
\begin{explain}
\begin{center}
\begin{tikzpicture}[scale=0.9,every node/.style={font=\small}]
\draw[->,gray] (-0.3,-0.3)--(-0.3,0.5) node[above]{N};
\fill (0,0) circle(2pt) node[below left]{Z};
\fill (0,2) circle(2pt) node[above left]{Y};
\fill (2.5,2) circle(2pt) node[above right]{X};
\fill (2.5,0) circle(2pt) node[below right]{P};
\draw[thick,brand] (0,0)--(0,2)--(2.5,2); \draw[thick,brand] (0,0)--(2.5,0);
\draw[-Stealth,thick,accent,dashed] (2.5,2)--(2.5,0.1);
\end{tikzpicture}
\end{center}
Y is north of Z. X is east of Y. P is east of Z. Taking the standard layout, X and P are at the same east--west position, with X directly above P. So P is to the \textbf{South} of X.
\end{explain}
\end{question}

\begin{question}{47}
{Match the following four terms of Column A with that of Column B, correctly as per Pramana's reasoning.\par\smallskip\begin{tabularx}{\linewidth}{@{}>{\raggedright\arraybackslash}X>{\raggedright\arraybackslash}X@{}}\toprule
\textbf{Column A} & \textbf{Column B}\\\midrule
1. Pramana & a. The subject who knows\\
2. Pramata & b. The means of validity\\
3. Prameya & c. The resultant valid knowledge\\
4. Prama & d. The object of knowledge\\
\bottomrule\end{tabularx}\par\smallskip
The correct match is\par}
\optstwo{1-d, 2-a, 3-b, 4-c}{1-b, 2-a, 3-d, 4-c}{1-c, 2-a, 3-b, 4-d}{1-d, 2-c, 3-b, 4-a}
\answer{B}{1-b, 2-a, 3-d, 4-c}
\begin{explain}
In Nyaya epistemology:
\begin{itemize}
\item \textbf{Pramana}: the means of valid knowledge (b)
\item \textbf{Pramata}: the knower, the subject who knows (a)
\item \textbf{Prameya}: the object of knowledge (d)
\item \textbf{Prama}: the resulting valid knowledge (c)
\end{itemize}
The match is \textbf{1-b, 2-a, 3-d, 4-c}.
\end{explain}
\end{question}

\begin{question}{48}
{The ratio of the number of students studying in schools A, B and C is 6:8:7 respectively. If the number of students studying in each of the schools is increased by 20\%, 15\% and 20\% respectively, what will be the new ratio of the number of students in schools A, B and C ?\par}
\optsfour{20:19:18}{20:22:24}{21:22:23}{18:23:21}
\answer{D}{18:23:21}
\begin{explain}
New ratio $=6\times1.20 : 8\times1.15 : 7\times1.20=7.2:9.2:8.4=72:92:84=\mathbf{18:23:21}$.
\end{explain}
\end{question}

\begin{question}{49}
{The average of $1^{2}$, $2^{2}$, $3^{2}$, $4^{2}$, $5^{2}$, $6^{2}$ is\par}
\optsfour{14.69}{15.17}{16.21}{17.48}
\answer{B}{15.17}
\begin{explain}
\[\frac{1+4+9+16+25+36}{6}=\frac{91}{6}=\mathbf{15.17}\]
(Using $\sum n^2=\tfrac{n(n+1)(2n+1)}{6}=\tfrac{6\cdot7\cdot13}{6}=91$.)
\end{explain}
\end{question}

\begin{question}{50}
{The bacterial count in a certain culture (colony) increased from 1000 to 4000 in 3 days. What was the average percentage increase per day?\par}
\optsfour{49.2\%}{54.3\%}{58.7\%}{62.9\%}
\answer{C}{58.7\%}
\begin{explain}
The growth is compound (per day): $1000(1+r)^3=4000\Rightarrow(1+r)^3=4\Rightarrow1+r=\sqrt[3]4\approx1.587$.
Average daily increase $\approx\mathbf{58.7\%}$.
\end{explain}
\end{question}
